% 第 4 章 §4.2.1 后段–§4.2.3 开头（印刷页 172–178，PDF 182–188；原文 ch04.txt 第 332–795 行）
% 承接 ch04a：该文件止于引理 4.2.8 之前的引导句，本文件自引理 4.2.8 起，故不另写 \section。
% 定理类共享计数器：本块首个环境为引理 4.2.8，故置 theorem = 7。
\setcounter{theorem}{7}

\begin{lemma}\label{lem:4.2.8}
设 $f\colon(E_1,\Phi_1)\to(E_2,\Phi_2)$ 是 Higgs 丛的全纯同态，且 $\Phi_2f=f\Phi_1$。设 $(E_i,\Phi_i)$ 半稳定，$i=1,2$，而 $\mu(E_1)>\mu(E_2)$。则 $f\equiv0$。若 $\mu(E_1)=\mu(E_2)$，且两者之一是稳定的，则或者 $f\equiv0$，或者 $f$ 是同构。
\end{lemma}

\begin{proof}
先看第一个断言。此时若 $f\not\equiv0$，则由假设 $\Phi_2f=f\Phi_1$ 知 $f$ 的像是 $\Phi_2$-不变的；于是按斜率的条件，$f$ 必有核。但此时 $\ker f$ 是 $\Phi_1$-不变的。故 $\mu(\ker f)\le\mu(E_1)\le\mu(\operatorname{coker}f)\le\mu(E_2)$，矛盾。
\begin{annote}底本此处的不等式链中，$\mu(E_1)\le\mu(\operatorname{coker}f)$ 一步不能由半稳定性得到，且在 $\mu(E_1)>\mu(E_2)$ 之下一般不成立：$\operatorname{coker}f$ 是 $E_2$ 的商丛（带诱导 Higgs 场），半稳定性只给出 $\mu(\operatorname{coker}f)\le\mu(E_2)<\mu(E_1)$。译文照排底本的不等式链。按上下文，矛盾实际来自斜率在正合列上的可加性：$\mu(E_1)$ 居于 $\mu(\ker f)$ 与 $\mu(\operatorname{im}f)$ 之间，而 $\mu(\operatorname{im}f)\le\mu(E_2)<\mu(E_1)$，故 $\mu(\ker f)>\mu(E_1)$，与 $E_1$ 的半稳定性矛盾。[来源：上下文推断]\end{annote}
第二个断言类似可得。
\end{proof}

按定义，$(E,\Phi)$ 的\emph{Higgs 子丛}是 $\Phi$-不变的全纯子丛 $S\subset E$。$\Phi$ 在 $S$ 上的限制 $\Phi_S$ 使 $(S,\Phi_S)$ 成为 Higgs 丛，此时包含映射 $S\hookrightarrow E$ 给出 Higgs 丛的映射。类似地，$\Phi$ 在商丛 $Q=E/S$ 上诱导出 Higgs 丛结构。给定任意的 Higgs 丛，$(E,\Phi)$ 的 \emph{Harder--Narasimhan 滤过}（Harder--Narasimhan filtration，简称 HN 滤过）是一列由 Higgs 子丛构成的滤过
\[
0=(E_0,\Phi_0)\subset(E_1,\Phi_1)\subset\cdots\subset(E_\ell,\Phi_\ell)=(E,\Phi),
\]
使得各商丛 $(Q_i,\Phi_{Q_i})=(E_i,\Phi_i)/(E_{i-1},\Phi_{i-1})$ 半稳定（参看 [31]）。此外还要求该滤过满足 $\mu(Q_i)>\mu(Q_{i+1})$；可以证明，伴随分次对象（associated graded object）$\operatorname{Gr}_{\mathrm{HN}}(E,\Phi)=\bigoplus_{i=1}^{\ell}(Q_i,\Phi_{Q_i})$ 由 $(E,\Phi)$ 的同构类唯一确定。斜率之集 $\mu_i=\mu(Q_i)$ 是 Higgs 丛同构类的一个重要不变量。

\begin{remark}\label{rem:4.2.9}
由构造，$\mu_i$ 是 $E/E_{i-1}$ 连同其诱导 Higgs 场的 Higgs 子丛斜率的最大值。我们也可把 $\mu_i$ 解释为 $(E_i,\Phi_i)$ 的 Higgs 商丛斜率的最小值。事实上，$(E_1,\Phi_1)$ 半稳定，故 $i=1$ 时该断言平凡地成立。设 $i\ge2$，$E_i\to Q\to0$ 是满足 $\mu(Q)\le\mu_i$ 的 Higgs 商丛。若 $Q$ 是极小的这种商丛，则它关于诱导的 Higgs 场半稳定。由引理 4.2.8，诱导映射 $E_1\to Q$ 必为零。于是该商丛过渡为 $E/E_1\to Q\to0$。现在，由同样的论证，若 $i\ge3$，则 $E_2/E_1\to Q$ 为零。如此继续，我们得到商丛 $Q_i\to Q\to0$。又因 $(Q_i,\Phi_{Q_i})$ 半稳定而该商非零，再次应用引理 4.2.8，即得 $\mu_i\le\mu(Q)$。
\end{remark}

考虑由 HN 滤过得到的 $n$ 元数组 $\vec\mu(E,\Phi)=(\mu_1,\dots,\mu_n)$，其中每个 $\mu_i$ 按 $Q_i$ 的秩重复。于是得到向量 $\vec\mu(E,\Phi)$，称为 $(E,\Phi)$ 的 \emph{Harder--Narasimhan 型}（HN 型）。这类向量上有一个自然的偏序（partial ordering），它正是我们所需分层的要点。对满足 $\mu_1\ge\cdots\ge\mu_n$、$\lambda_1\ge\cdots\ge\lambda_n$ 且 $\sum_{i=1}^{n}\mu_i=\sum_{i=1}^{n}\lambda_i$ 的两个 $n$ 元数组 $\vec\mu,\vec\lambda$，我们定义
\[
\vec\lambda\le\vec\mu\quad\Longleftrightarrow\quad\sum_{j\le k}\lambda_j\le\sum_{j\le k}\mu_j\qquad\text{对一切 }k=1,\dots,n.
\]
这个序的重要性在于：它给出 Higgs 丛空间的一个分层（stratification）。特别地，Harder--Narasimhan 型是上半连续的（upper semicontinuous）。这正是全纯丛的 Atiyah--Bott 分层的直接类比 [2, §7]。

半稳定的 Higgs 丛 $(E,\Phi)$ 也有类似的滤过，其逐次商都稳定，且斜率都等于 $\mu(E)$。这称为 \emph{Seshadri 滤过}，见 [54]；其伴随分次对象 $\operatorname{Gr}_{\mathrm{S}}(E,\Phi)$ 从而是多稳定的。当 $\Phi\equiv0$ 时，我们恢复全纯丛 $E$ 通常的 Harder--Narasimhan 滤过与 Seshadri 滤过。我们将它们记作 $\operatorname{Gr}_{\mathrm{HN}}(E)$ 与 $\operatorname{Gr}_{\mathrm{S}}(E)$。

\begin{example}\label{ex:4.2.10}
考虑扩张 (4.5)，其中 $\operatorname{rank}S=\operatorname{rank}Q=1$ 且 $\deg S>\deg Q$。则 $E$ 的 Harder--Narasimhan 滤过由 $0\subset S\subset E$ 给出。
\end{example}

\subsection{模空间}

\paragraph{规范变换} 设 $\mathcal{A}_E$ 表示秩 $n$ 埃尔米特向量丛 $E$ 上酉联络的空间。若 $\mathfrak{g}_E$ 表示 $E$ 的伴随反埃尔米特自同态丛，则由 Leibniz 法则可以看出，$\mathcal{A}_E$ 是以 $\Omega^1(X,\mathfrak{g}_E)$ 为模型的无穷维仿射空间。由 §4.2.1 中讨论的Chern 连接的构造还可看出，$\mathcal{A}_E$ 可以与 $E$ 上全纯结构的空间等同。我们最常关心的是固定行列式的情形，即 $\det E$ 上诱导的全纯结构固定。

\emph{规范群}定义为 $\mathcal{G}_E=\{g\in\Omega^0(X,\operatorname{End}E): gg^*=I\}$。在固定行列式的情形，我们还加上条件 $\det g=1$（见注 4.2.1）。规范群通过拉回联络作用于 $\mathcal{A}_E$：$d_{g(A)}=g\circ d_A\circ g^{-1}$。另一方面，由于与全纯结构的等同，复化 $\mathcal{G}_E^\C$ 即\emph{复规范群}也作用于 $\mathcal{A}_E$。具体地说，若 $\bar\partial_E=d''_A$，则 $g(A)$ 是 $g\circ\bar\partial_E\circ g^{-1}$ 的Chern 连接。

Higgs 丛的空间是
\[
\mathcal{B}_E=\{(A,\Phi)\in\mathcal{A}_E\times\Omega^0(X,K\otimes\mathfrak{g}^\C_E): d''_A\Phi=0\}.
\]
设 $\mathcal{B}^{ss}_E\subset\mathcal{B}_E$ 表示半稳定 Higgs 丛的子集。

\begin{definition}\label{def:4.2.11}
$X$ 上秩 $n$ 半稳定 Higgs 丛（带固定行列式）的模空间是 $\mathfrak{M}_D^{(n)}=\mathcal{B}^{ss}_E\sslash \mathcal{G}_E^\C$，其中双斜线表示把 $(E,\Phi)$ 与 $\operatorname{Gr}_{\mathrm{S}}(E,\Phi)$ 的轨道等同起来。
\end{definition}

我们一直没有仔细处理拓扑。事实上，利用 Kuranishi 映射（参看 [44]），可以给 $\mathfrak{M}_D^{(n)}$ 赋予（可能非既约的）复解析空间（complex analytic space）的结构。用几何不变量理论的代数构造见 [60]。

第二点说明是：$\mathcal{G}_E^\C/\mathcal{G}_E$ 可以与 $E$ 上埃尔米特度量的空间等同。在研究泛函沿 $\mathcal{A}_E/\mathcal{G}_E$ 中 $\mathcal{G}_E^\C$ 轨道的行为时，这导致一个重要的解释：我们既可以认为是在固定的埃尔米特度量下改变复结构 $g(\bar\partial_E)$，也可以固定 $\bar\partial_E$ 而按 $\langle s_1,s_2\rangle_{g(H)}=\langle gs_1,gs_2\rangle_H$ 改变度量 $H$。

\paragraph{Higgs 丛的形变} 令 $D''=d''_A+\Phi$，$D'=d'_A+\Phi^*$。$X$ 上的度量 $\omega$ 与 $E$ 上的埃尔米特度量在取值于 $E$ 与 $\operatorname{End}E$ 的形式上定义 $L^2$ 内积。我们有\emph{凯勒恒等式}（Kähler identities）
\begin{equation}
\begin{aligned}
(D'')^*&=-\sqrt{-1}[\Lambda,D'],\\
(D')^*&=\sqrt{-1}[\Lambda,D''],
\end{aligned}
\label{eq:4.6}\tag{4.6}
\end{equation}
（$\Phi=0$ 的情形见 [26, p.~111]；$\Phi\ne0$ 的情形可由直接计算得到）。

模空间的无穷小结构由一个\emph{形变复形}（deformation complex）$C(A,\Phi)$ 支配，后者通过对条件 $d''_A\Phi=0$ 与规范群的作用求导得到：
\begin{equation}
\begin{aligned}
C(A,\Phi)\colon 0\longrightarrow\Omega^0(X,\mathfrak{g}^\C_E)&\xrightarrow{\,D''\,}\Omega^{1,0}(X,\mathfrak{g}^\C_E)\oplus\Omega^{0,1}(X,\mathfrak{g}^\C_E)\\
&\xrightarrow{\,D''\,}\Omega^{1,1}(X,\mathfrak{g}^\C_E)\longrightarrow0.
\end{aligned}
\label{eq:4.7}\tag{4.7}
\end{equation}
注意 $\Phi$ 的全纯性条件保证 $(D'')^2=0$。Serre 对偶给出同构 $H^0(C(A,\Phi))\simeq H^2(C(A,\Phi))$。若 $H^0(C(A,\Phi))\simeq\C$（固定行列式情形为 $\{0\}$），则称该 Higgs 丛是\emph{单的}（simple）。

\begin{remark}\label{rem:4.2.12}
稳定的 Higgs 丛必是单的。事实上，若 $\varphi\in\ker D''$，则 $\varphi$ 是 $E$ 的与 $\Phi$ 交换的全纯自同态。特别地，$\det\varphi$ 是全纯函数，因而是常数。又 $\ker\varphi$ 是 $\Phi$-不变的。若 $\varphi$ 非零但不是同构，则
\[
0\longrightarrow\ker\varphi\longrightarrow E\longrightarrow E/\ker\varphi\longrightarrow0.
\]
由于 $E/\ker\varphi$ 也是 $E$ 的子层，稳定性蕴涵 $\mu(\ker\varphi)$ 与 $\mu(E/\ker\varphi)$ 都小于 $\mu(E)$，矛盾。故 $\varphi$ 或为零或是同构。但对任意标量 $\lambda$ 把同样的论证用于 $\varphi-\lambda$，即知 $\varphi$ 是单位矩阵的倍数。
\end{remark}

\begin{proposition}\label{prop:4.2.13}
在单 Higgs 丛 $[A,\Phi]$ 处，$\mathfrak{M}_D^{(n)}$ 是光滑的，复维数为 $(n^2-1)(2g-2)$，而切空间可等同于
\begin{equation}
H^1(C(A,\Phi))\simeq\Bigl\{(\varphi,\beta): d''_A\varphi=-[\Phi,\beta],\ (d''_A)^*\beta=\sqrt{-1}\Lambda[\Phi^*,\varphi]\Bigr\}.
\label{eq:4.8}\tag{4.8}
\end{equation}
\end{proposition}

\begin{example}[参看 [35, 38]]\label{ex:4.2.14}
现在给出稳定 Higgs 丛的重要例子，即 Fuchsian 的那些。先看秩 $2$ 的情形。选定典范丛的一个平方根（square root）$K^{1/2}$，并令 $E=K^{1/2}\oplus K^{-1/2}$。于是自同态丛中把 $K^{1/2}$ 送到 $K^{-1/2}$ 的那一部分同构于 $K^{-1}$；张量上 $K$ 之后，它变为平凡丛。因此
\[
\Phi=\begin{pmatrix}0&1\\0&0\end{pmatrix}
\]
作为 Higgs 场是有意义的，且它显然是全纯的。虽然 $E$ 作为全纯向量丛是不稳定的，但 Higgs 丛 $(E,\Phi)$ 稳定，因为唯一的 $\Phi$-不变线子丛是 $K^{-1/2}$，而它的次数为负。顺带说明：若考虑 $E$ 上由 $\bar\partial$-算子
\[
\bar\partial_E+\Phi^*=\begin{pmatrix}\bar\partial_{K^{1/2}}&\omega\\0&\bar\partial_{K^{-1/2}}\end{pmatrix}
\]
给出的另一个全纯结构 $V$，则 $V$ 是唯一（同构意义下）的非分裂扩张
\[
0\longrightarrow K^{1/2}\longrightarrow V\longrightarrow K^{-1/2}\longrightarrow0.
\]
现在来计算 $[(E,\Phi)]$ 处的切空间 $\mathfrak{M}_D^{(2)}$。记
\[
\beta=\begin{pmatrix}b&b_1\\b_2&-b\end{pmatrix},\qquad\varphi=\begin{pmatrix}\varphi&\varphi_1\\\varphi_2&-\varphi\end{pmatrix},
\]
并计算
\[
[\Phi,\beta]=\begin{pmatrix}-b_1&0\\2b&b_1\end{pmatrix},\qquad\sqrt{-1}\Lambda[\Phi^*,\varphi]=\begin{pmatrix}\varphi_2&-2\varphi\\0&-\varphi_2\end{pmatrix}.
\]
于是 $(\beta,\varphi)$ 定义切向量的条件 (4.8) 为
\[
\bar\partial_E\varphi=\begin{pmatrix}b_1&0\\-2b&-b_1\end{pmatrix},\qquad\bar\partial^*_E\beta=\begin{pmatrix}\varphi_2&-2\varphi\\0&-\varphi_2\end{pmatrix}.
\]
特别地，$\varphi_1\in H^0(X,K^2)$，$b_2\in H^{0,1}_{\bar\partial}(X,K^*)\simeq H^0(X,K^2)^*$。断言其余分量全为零。事实上，$\varphi$ 与 $b_1$ 所满足的方程是 $\bar\partial\varphi=b_1$ 与 $\bar\partial^*b_1=-2\varphi$。但这蕴涵 $(\bar\partial^*\bar\partial+2)\varphi=0$。故 $\varphi$ 从而 $b_1$ 必为零。对 $\varphi_2$ 与 $b$ 有同样的论证。因此得到同构
\[
T_{[E_F,\Phi_F]}\mathfrak{M}_D^{(2)}\simeq H^0(X,K^2)\oplus(H^0(X,K^2))^*.
\]
对 $n\ge2$ 有类似的论证。这里取
\[
E_F=K^{(n-1)/2}\oplus K^{(n-3)/2}\oplus\cdots\oplus K^{-(n-1)/2}
\]
以及
\[
\Phi_F=\begin{pmatrix}
0&0&0&\cdots&0\\
1&0&0&\cdots&0\\
0&1&0&\cdots&0\\
\vdots& &\ddots&\ddots&\vdots\\
0&\cdots&0&1&0
\end{pmatrix}.
\]
注意，关于该分裂，$\varphi$ 的 $(i,j)$ 元是 $K^{j-i+1}$ 的截面，而 $\beta$ 的 $(i,j)$ 元属于 $\Omega^{0,1}(X,K^{j-i})$。对切向量 $(\beta,\varphi)$ 的各分量得到如下方程：
\begin{equation}
\begin{aligned}
\bar\partial_E\varphi_{i,j}&=\beta_{i-1,j}-\beta_{i,j+1},\\
\bar\partial^*_E\beta_{i,j}&=\varphi_{i,j-1}-\varphi_{i+1,j},
\end{aligned}
\label{eq:4.9}\tag{4.9}
\end{equation}
其中约定下标 $\le0$ 或 $\ge n+1$ 的项取为零。与 $n=2$ 的情形一样再求导一次，得到
\begin{equation}
\begin{aligned}
(L-\delta_{i,1}-\delta_{j,n})\varphi_{i,j}&=\varphi_{i+1,j+1}+\varphi_{i-1,j-1},\\
(\widetilde{L}-\delta_{i,n}-\delta_{j,1})\beta_{i,j}&=\beta_{i+1,j+1}+\beta_{i-1,j-1},
\end{aligned}
\label{eq:4.10}\tag{4.10}
\end{equation}
其中 $L=\bar\partial^*_E\bar\partial_E+2$，$\widetilde{L}=\bar\partial_E\bar\partial^*_E+2$。断言：当 $i\ge j$ 时 $\varphi_{i,j}=0$（相应地，当 $i\le j$ 时 $\beta_{i,j}=0$）。例如，由 (4.10) 有 $L\varphi_{n,1}=0$，而 $L$ 是正算子，故 $\varphi_{n,1}$ 为零。更一般地，固定 $0\le p\le n-2$。则对 $0\le\ell\le n-p-1$，存在多项式 $P_\ell$ 使得
\begin{equation}
\varphi_{p+\ell+1,\ell+1}=P_\ell(L)\varphi_{p+1,1}.
\label{eq:4.11}\tag{4.11}
\end{equation}
事实上，令 $P_0(L)=1$；若 $p\ne0$ 令 $P_1(L)=L$，若 $p=0$ 令 $P_1(L)=L-1$。设 $P_k(L)$ 已对 $0\le k\le\ell$ 定义，其中 $0<\ell<n-p-1$。利用 (4.10) 与 (4.11) 得到
\[
\begin{aligned}
&L\varphi_{p+\ell+1,\ell+1}=\varphi_{p+\ell+2,\ell+2}+\varphi_{p+\ell,\ell},\\
&LP_\ell(L)\varphi_{p+1,1}=\varphi_{p+\ell+2,\ell+2}+P_{\ell-1}(L)\varphi_{p+1,1}.
\end{aligned}
\]
于是令 $P_{\ell+1}(L)=LP_\ell(L)-P_{\ell-1}(L)$。由于 $L\ge2$，由递归定义看出 $P_{\ell+1}(L)\ge P_\ell(L)$，从而对一切 $\ell\ge1$ 有 $P_\ell(L)\ge P_1(L)\ge1$，且当 $p\ne0$ 时 $\ge2$。在 (4.11) 中取 $\ell=n-p-1$，得
\begin{equation}
\varphi_{n,n-p}=P_{n-p-1}(L)\varphi_{p+1,1}.
\label{eq:4.12}\tag{4.12}
\end{equation}
另一方面，类似的论证蕴涵 $\varphi_{p+1,1}=P_{n-p-1}(L)\varphi_{n,n-p}$，由此得到
\[
0=(P^2_{n-p-1}(L)-1)\varphi_{p+1,1}=(P_{n-p-1}(L)+1)(P_{n-p-1}(L)-1)\varphi_{p+1,1}.
\]
于是 $\varphi_{p+1,1}$ 属于 $P_{n-p-1}(L)-1$ 的核。但由上面的说明，当 $p\ge1$ 时 $\varphi_{p+1,1}$ 必为零。由于 $p\ge1$ 是任意的，这由 (4.11) 蕴涵 $i>j$ 时 $\varphi_{i,j}=0$。在 $p=0$ 的情形，注意对一切 $\ell\ge1$，$P_\ell(L)$ 是 $\bar\partial^*_E\bar\partial_E$ 的正次数多项式，其系数非负且常数项等于 $1$。事实上，由定义
\[
P_{\ell+1}(L)-P_\ell(L)=(\bar\partial^*_E\bar\partial_E)P_\ell(L)+P_\ell(L)-P_{\ell-1}(L),
\]
故由归纳法，$P_{\ell+1}(L)-P_\ell(L)$ 的系数非负且常数项为零。此时 $(P_{n-1}(L)-1)\varphi_{1,1}=0$ 蕴涵 $\varphi_{1,1}$ 全纯。再次利用 (4.11)，
\[
\varphi_{\ell+1,\ell+1}=P_\ell(L)\varphi_{1,1}=(P_\ell(L)-1)\varphi_{1,1}+\varphi_{1,1}=\varphi_{1,1}
\]
对一切 $\ell=0,\dots,n-1$ 成立。但由于 $(\varphi_{i,j})$ 无迹，实际上对一切 $i$ 有 $\varphi_{i,i}=0$。对 $\beta_{i,j}$ 的证明类似。

回到 (4.9)，我们看到：当 $i<j$ 时，$\varphi_{i,j}$ 全纯（相应地，$\beta_{j,i}$ 调和）。此外，对 $p\ge1$，(4.10) 变为
\begin{equation}
(2-\delta_{i,1}-\delta_{i,n-p})\varphi_{i,i+p}=\varphi_{i+1,i+1+p}+\varphi_{i-1,i-1+p}.
\label{eq:4.13}\tag{4.13}
\end{equation}
若 $i=1$，这蕴涵 $\varphi_{1,p+1}=\varphi_{2,p+2}$。归纳假设对一切 $k\le i$ 有 $\varphi_{k,k+p}=\varphi_{1,p+1}$。则若 $i+p\ne n$，(4.13) 蕴涵
\[
2\varphi_{i,i+p}=\varphi_{i+1,i+1+p}+\varphi_{i-1,i-1+p}\quad\Longrightarrow\quad\varphi_{1,p+1}=\varphi_{i+1,i+1+p}.
\]
若 $i+p=n$，则立即得到 $\varphi_{i,n}=\varphi_{i-1,n-1}=\varphi_{1,p+1}$。因此，所有满足 $j-i=p$ 的 $\varphi_{i,j}$ 都相等。对 $\beta_{i,j}$ 有同样的论证。由此断定，映射 $(\varphi,\beta)\mapsto(\varphi_{1,2},\dots,\varphi_{1,n},\beta_{2,1},\dots,\beta_{n,1})$ 给出同构
\begin{equation}
T_{[E_F,\Phi_F]}\mathfrak{M}_D^{(n)}\simeq\bigoplus_{j=2}^{n}H^0(X,K^j)\oplus(H^0(X,K^j))^*.
\label{eq:4.14}\tag{4.14}
\end{equation}
以 $\bar\partial_E+\Phi^*_F$ 为 $\bar\partial$-算子的秩 $n$ 全纯向量丛 $V$ 不稳定，并有 Harder--Narasimhan 滤过 $0=V_0\subset V_1\subset\cdots\subset V_n=V$，$V_{j+1}/V_j=K^{-j+(n-1)/2}$，使得
\[
0\longrightarrow V_j\longrightarrow V_{j+1}\longrightarrow K^{-j+(n-1)/2}\longrightarrow0
\]
是（唯一的）非分裂扩张。这是 oper 的一个例子。oper 将在 §4.4.3 中讨论。
\end{example}

\begin{annote}底本 §4.2.2 交替使用正体 $K$ 与花体 $\mathcal{K}$ 表示典范丛：p.~173 的 $\mathcal{B}_E$ 定义作 $K$，p.~174 起多作 $\mathcal{K}$，而 p.~175 同一行内先作 $H^{0,1}_{\bar\partial}(X,K^*)$、后作 $H^0(X,\mathcal{K}^2)^*$。译文依 ch04a 的记号表（$K=K_X$ 为典范丛）一律作 $K$，不逐处标注。[来源：上下文推断]\end{annote}

\paragraph{Hitchin 映射} 给定 Higgs 丛 $(E,\Phi)$，$\det(\lambda+\Phi)$ 的展开式中 $\lambda^{n-i}$ 的系数是 $K^i$ 的全纯截面，$i=1,\dots,n$。在我们主要考虑的固定行列式情形，$\operatorname{tr}\Phi=0$，故这些截面从 $i=2$ 开始。这些多重典范截面（pluricanonical sections）显然在 $\mathcal{G}_E^\C$ 的（共轭）作用下不变，于是我们得到一个良定义的映射，称为 \emph{Hitchin 映射}，
\begin{equation}
h\colon\mathfrak{M}_D^{(n)}\longrightarrow\bigoplus_{i=2}^{n}H^0(X,K^i).
\label{eq:4.15}\tag{4.15}
\end{equation}
该映射的结构及其纤维极为丰富（参看 [36]）。不过，在这些笔记中我只讨论下面这个重要事实，它将在下一节用 Uhlenbeck 紧性证明（代数证明见 [50, 59]）。\begin{annote}底本此句印作 the fibers turns out be extremely rich，脱落 to（应为 turns out to be）；中文不体现该处笔误，译文按正字处理。[来源：上下文推断]\end{annote}

\begin{theorem}\label{thm:4.2.15}
Hitchin 映射是紧合的。
\end{theorem}

\subsection{Hitchin--Kobayashi 对应}

\paragraph{稳定性与临界度量} 平凡丛上 Higgs 丛的 \emph{Hitchin 方程}是
\begin{equation}
F_A+[\Phi,\Phi^*]=0.
\label{eq:4.16}\tag{4.16}
\end{equation}
这里 $\Phi$ 被视为取值于自同态的 $(1,0)$-形式。考虑次数非零的丛也是方便的。此时方程变为
\begin{equation}
f_{(A,\Phi)}:=\sqrt{-1}\Lambda(F_A+[\Phi,\Phi^*])=\mu.
\label{eq:4.17}\tag{4.17}
\end{equation}
这里回忆归一化 $\operatorname{vol}(X)=2\pi$；于是右端单位自同态的标量倍数必然满足 $\mu=\mu(E)$。

对 (4.17) 有两种思考方式：对 Higgs 丛 $(E,\Phi)$，埃尔米特度量的选取给出Chern 连接 $A=(\bar\partial_E,H)$。因此，我们或者把 (4.17) 视为关于埃尔米特度量 $H$ 的方程，或者（等价地）固定 $H$ 而考虑复规范轨道中所有 $(A,\Phi)$ 的 $f_{(A,\Phi)}$。我们将经常在这两种等价的看法之间来回。

方程 (4.17) 的解可视为全纯对空间上 \emph{Yang--Mills--Higgs 泛函}的绝对极小（absolute minimum），该泛函定义为
\begin{equation}
\operatorname{YMH}(A,\Phi)=\int_X|F_A+[\Phi,\Phi^*]|^2\omega.
\label{eq:4.18}\tag{4.18}
\end{equation}
YMH 的 Euler--Lagrange 方程是
\begin{equation}
d_Af_{(A,\Phi)}=0,\qquad[\Phi,f_{(A,\Phi)}]=0.
\label{eq:4.19}\tag{4.19}
\end{equation}
若 (4.19) 成立，则称度量为\emph{临界}的。此时容易看出，丛 $(E,\Phi)$ 全纯且等距地分裂为若干 Higgs 丛的直和，这些 Higgs 丛是 (4.17) 的解，但斜率可以不同。

\begin{proposition}\label{prop:4.2.16}
若 Higgs 丛 $(E,\Phi)$ 容许满足 (4.17) 的度量，则 $(E,\Phi)$ 是多稳定的。
\end{proposition}

\begin{proof}
设 $S\subset E$ 是真 $\Phi$-不变子丛。令 $\pi$ 表示到 $S$ 的正交投影，$\beta=-\bar\partial_E\pi$ 是第二基本形式。则由于 $S$ 不变，$(I-\pi)\Phi\pi=0$，即 $\Phi\pi=\pi\Phi\pi$、$\pi\Phi^*=\pi\Phi^*\pi$。特别地，这蕴涵
\begin{align*}
\operatorname{tr}\bigl(\pi[\Phi,\Phi^*]\bigr)&=\operatorname{tr}(\pi\Phi\Phi^*)+\operatorname{tr}(\pi\Phi^*\Phi)\\
&=\operatorname{tr}(\pi\Phi\Phi^*)-\operatorname{tr}(\Phi\pi\Phi^*)\\
&=\operatorname{tr}(\pi\Phi\Phi^*\pi)-\operatorname{tr}(\Phi\pi\Phi^*\pi)\\
&=\operatorname{tr}(\pi\Phi\Phi^*\pi)-\operatorname{tr}(\pi\Phi\pi\Phi^*\pi)\\
&=\operatorname{tr}\bigl(\pi\Phi(I-\pi)\Phi^*\pi\bigr)=\operatorname{tr}\bigl(\pi\Phi(I-\pi)(I-\pi)\Phi^*\pi\bigr)\\
&=\operatorname{tr}\bigl\{(\pi\Phi(I-\pi))(\pi\Phi(I-\pi))^*\bigr\},
\end{align*}
\end{proof}

% 接缝：命题 4.2.16 的证明在 p.178 页末中断，余下部分（原文「and tr(π√−1Λ[Φ,Φ*])=…」(4.20) 一节）
% ch04c 应紧接一句「and tr(π√−1Λ[Φ,Φ*]) = …」(4.20)，并以「命题 4.2.16 的证明（续）」为 proof 环境标题续写（体例同 ch04e.tex）。
